%==============================================================================
% PAPER B — APPLICATION PAPER
% A control case study made easy by CosimGym (incremental: B1 -> B2 -> B3).
% Target: Energy and Buildings (B1) or Applied Energy (B1+B2).
% STATUS: first prose draft. Intro/Related/Methodology = real draft text.
% Results section = guidelines; numbers/figures added ONLY after experiments [TODO].
% Citations are \cite{...} placeholders; share cosimgym.bib with Paper A.
%==============================================================================
\documentclass[preprint,review,12pt]{elsarticle}
\usepackage{graphicx,booktabs,amsmath,siunitx,hyperref}

\journal{Energy and Buildings} % escalate to Applied Energy if B2 included

\begin{document}
\begin{frontmatter}

\title{Multi-rate Reinforcement Learning Control of an EnergyPlus-Coupled\\
Heat-Pump Building via Declarative Co-Simulation}

\author[polito]{Pietro Rando Mazzarino\corref{cor}}
\ead{pietro.randomazzarino@polito.it}
\author[polito]{Co-authors}
\address[polito]{EC-Lab, Politecnico di Torino, Turin, Italy}
\cortext[cor]{Corresponding author.}

\begin{abstract}
Heat pumps are central to building decarbonization and, as flexible electric
loads, to demand response, which motivates intelligent control able to balance
thermal comfort against energy use and cost. Reinforcement learning (RL) is a
promising approach, but its reported performance depends heavily on the simulation
used for training and evaluation, and most studies rely on a single fixed-rate
simulator with bespoke, hard-to-reproduce integration. We study RL control of a
heat-pump-heated building modeled in EnergyPlus and coupled, at their natural and
distinct time rates, to a heat-pump model and a weather feed, using the CosimGym
co-simulation framework so that the entire multi-rate, multi-tool environment is
assembled purely from configuration. We train DQN and SAC agents and compare them
against rule-based, PID, and model-predictive control baselines on energy, thermal
comfort, and cost.
% TODO: 1-2 sentences with headline numbers after experiments.
We further analyze how the episode-reset strategy required by a non-resettable
co-simulation affects learning, and we assess generalization across climates. The
study shows that realistic, multi-formalism physics can be made a configuration
concern, enabling fair and reproducible comparison between RL and model-based
control.
\end{abstract}

% NOTE: elsarticle has NO \begin{highlights} environment (Elsevier CAS class only).
% Submit Highlights as a separate file. Draft highlights:
%  - EnergyPlus-FMU building, heat pump and weather coupled at native multi-rates
%  - DQN/SAC trained purely from configuration; compared against PID, MPC, rule-based
%  - Episode-reset strategy shown to affect learning over a non-resettable co-simulation
%  - Cross-climate generalization of the learned policy
%  - Whole multi-rate, multi-formalism study assembled without orchestration code

\begin{keyword}
reinforcement learning \sep building energy management \sep heat pump \sep
co-simulation \sep EnergyPlus \sep model predictive control \sep demand response
\end{keyword}
\end{frontmatter}

%==============================================================================
\section{Introduction}\label{sec:intro}
%==============================================================================
Buildings account for a large share of global final energy use and
emissions~\cite{iea_buildings}, and space heating is among the largest end uses.
The electrification of heating through heat pumps is a key decarbonization lever,
and because heat pumps are deferrable, modulating electric loads coupled to
buildings with significant thermal inertia, they are also a valuable source of
demand-side flexibility in increasingly renewable-dominated power
systems~\cite{rl_energy_review}. Realizing this potential requires control that
goes beyond fixed schedules and simple thermostats, jointly managing occupant
comfort, energy consumption, and operating cost under time-varying weather, prices,
and occupancy.

Reinforcement learning has attracted strong interest for this problem because it
can learn adaptive policies directly from interaction, without an explicit
optimization model at run time~\cite{drl_hvac_review}. However, RL agents are
trained and evaluated in simulation, and the credibility of reported results is
therefore inseparable from the realism and rigor of that simulation. Two
limitations recur in the literature. First, most studies use a single simulator at
a single, fixed time resolution, even though the relevant physics spans very
different time scales---weather and thermal mass evolve over hours, while the heat
pump and its controller act over minutes. Second, the coupling between the
simulator and the RL stack is typically a one-off engineering effort, which makes
results hard to reproduce and hard to compare fairly against established baselines
such as model-predictive control (MPC)~\cite{boptest}.

A specific consequence is that coupling a high-fidelity building model, such as an
EnergyPlus model, with a faster heat-pump and controller loop and a slower weather
process \emph{at their natural rates}, and then comparing RL against MPC under
exactly the same physics, is rarely done---not because it is uninteresting, but
because the required co-simulation plumbing is laborious. Co-simulation middleware
such as HELICS~\cite{helics} and the FMI standard~\cite{fmi} make such coupling
technically possible, yet wiring them to an RL training loop has, until now,
remained a bespoke task.

This paper removes that obstacle and exploits the result. Using the CosimGym
framework~\cite{cosimgym_paperA}, which compiles a declarative scenario into a
HELICS co-simulation exposed as a standard Gymnasium environment, we assemble a
multi-rate, EnergyPlus-coupled heat-pump control problem entirely from
configuration. Our contributions are: (i) a multi-rate, multi-formalism RL control
study of a heat-pump building in which the realistic environment is a configuration
artifact rather than custom code; (ii) a quantitative comparison of DQN and SAC
against rule-based, PID, and MPC baselines on energy, comfort, and cost; (iii) an
analysis of how the episode-reset strategy, made necessary by a non-resettable
co-simulation, affects learning; and (iv) an assessment of cross-climate
generalization. Section~\ref{sec:related} reviews related work,
Section~\ref{sec:method} describes the environment and methods,
Section~\ref{sec:results} reports results, and
Sections~\ref{sec:future}--\ref{sec:conclusion} discuss limitations and conclude.

%==============================================================================
\section{Related work}\label{sec:related}
%==============================================================================
\paragraph{RL for building and HVAC control}
Deep RL has been applied extensively to building climate and HVAC control, with
reported gains in energy and comfort across a range of building types and
algorithms~\cite{drl_hvac_review}. A recurring theme in critical reviews is that
comparability is hampered by heterogeneous, simulator-specific experimental setups,
which motivates standardized environments and protocols.

\paragraph{RL environments and co-simulation for buildings}
Standardized testbeds have emerged to address this, including
Sinergym~\cite{sinergym} over EnergyPlus, BOPTEST~\cite{boptest} over Modelica
emulators with explicit support for benchmarking RL against MPC, Energym over
FMUs~\cite{energym}, and CityLearn~\cite{citylearn} for district-level
coordination. These are valuable but are each tied to a particular simulator or
domain and to a fixed temporal structure. General co-simulation middleware such as
HELICS~\cite{helics} and Mosaik~\cite{mosaik}, together with FMI~\cite{fmi},
provides the coupling and interoperability we rely on, but offers no native RL
interface.

\paragraph{RL versus model-predictive control}
MPC is the established advanced-control baseline for buildings, and fair RL-versus-
MPC comparisons under identical physics are an active concern~\cite{boptest}. By
making the physics a shared configuration artifact, our setup allows both
controllers to act on exactly the same multi-rate co-simulation. To our knowledge,
no prior building-control study combines native multi-rate coupling, multiple model
formalisms, configuration-only composition, and reset-aware episodic training, all
of which we exploit here.

%==============================================================================
\section{Methodology}\label{sec:method}
%==============================================================================
\subsection{Co-simulation environment}
We build the environment on CosimGym~\cite{cosimgym_paperA}, which compiles a
single declarative scenario into a HELICS federation that is simultaneously a
Gymnasium environment; we summarize only what is needed here and refer the reader
to the framework paper for details. The federation, shown in
Figure~\ref{fig:federation}, comprises a weather federate, a building federate, a
heat-pump federate, and an RL agent. Crucially, the federates run at different real
periods---weather updates hourly, the building model at a sub-hourly step, and the
heat pump and controller at a one-minute step---so that the agent learns against
genuinely multi-rate dynamics. The complete environment, including data flow,
timing, and the agent's wiring, is specified in one YAML file (provided in the
appendix and repository), and all runs are logged to a time-series store for
traceability.

\begin{figure}[t]\centering
% TODO: data-flow diagram of the federation with per-federate rates.
\fbox{\parbox{0.9\linewidth}{\centering\vspace{2.2cm}
\textbf{[Figure 1: federation data flow + rates]}\vspace{2.2cm}}}
\caption{Multi-rate federation: weather (hourly) drives the EnergyPlus building
model and the heat pump; the building temperature and external conditions form the
agent's observations; the agent's modulation command actuates the heat pump.}
\label{fig:federation}
\end{figure}

\subsection{Physical models}
The building is represented by an EnergyPlus model exported as an FMI~2.0 FMU,
capturing envelope, internal gains, and zone thermal response; the heat pump is
modeled with a temperature-dependent coefficient of performance; and weather is
streamed from measured data. Model parameters are listed in
Table~\ref{tab:params}. We sanity-check the building model under free-floating and
step-input conditions before control experiments.
% TODO: Table 1 with concrete parameters; brief validation note + figure if useful.

\begin{table}[t]\centering
\caption{Key model parameters.}
\label{tab:params}
\small
\begin{tabular}{@{}lll@{}}
\toprule
Component & Parameter & Value \\
\midrule
Building (EnergyPlus FMU) & \textit{[TODO]} & \textit{[TODO]} \\
Heat pump & rated power, COP model & \textit{[TODO]} \\
Weather & location, year & \textit{[TODO]} \\
\bottomrule
\end{tabular}
\end{table}

\subsection{Control problem formulation}
The agent observes the indoor temperature and external conditions (optionally
augmented with time-of-day or price signals) and outputs the heat-pump modulation
command. The reward at each step penalizes thermal discomfort and energy use,
\begin{equation}
r_t = -\,w_c\,\big[d(T^{\text{in}}_t, \mathcal{C})\big]
      -\,w_e\,E_t ,
\label{eq:reward}
\end{equation}
where $d(\cdot,\mathcal{C})$ measures deviation of the indoor temperature from the
comfort band $\mathcal{C}$, $E_t$ is the electrical energy consumed in the step,
and $w_c, w_e$ weight the objectives (a cost term $w_p\,c_t E_t$ under a tariff is
added for the prosumer extension). Episodes are defined over the advancing
federated clock; because the co-simulation cannot be rewound freely, we use the
reset strategy described next.

\subsection{Agents, baselines, and reset strategy}
We train a discrete-action DQN and a continuous-action SAC agent (Stable-Baselines3
implementations) and compare against three baselines: a rule-based controller, a
tuned PID controller, and an MPC controller using a reduced building model over a
finite horizon. Hyperparameters are reported in Table~\ref{tab:hparams}. Because
the federation is not freely resettable, training uses the framework's reset
mechanism; we adopt a rolling reset as the default and treat the reset strategy
itself as an experimental factor in Section~\ref{sec:results}.
% TODO: Table 2 hyperparameters; MPC model + horizon; seeds N.

\begin{table}[t]\centering
\caption{Agent and baseline configuration.}
\label{tab:hparams}
\small
\begin{tabular}{@{}lll@{}}
\toprule
Controller & Key settings & Value \\
\midrule
DQN / SAC & learning rate, $\gamma$, batch, steps & \textit{[TODO]} \\
PID       & gains, setpoint & \textit{[TODO]} \\
MPC       & model, horizon, solver & \textit{[TODO]} \\
\bottomrule
\end{tabular}
\end{table}

\subsection{Evaluation protocol}
All controllers are evaluated on a held-out period at the same conditions. We
report heating electrical energy, a thermal-comfort violation metric integrated
over time, operating cost under a representative tariff, and heat-pump efficiency
utilization, as mean and standard deviation over multiple random seeds. We provide
seeds, software versions, and configuration hashes to support reproduction.

%==============================================================================
\section{Results and discussion}\label{sec:results}
%==============================================================================
% Guidelines only until experiments are run. Do NOT fill numbers before then.

\subsection{Control performance versus baselines}
% Lead with the headline. Fig. 2: temperature and power trajectories for
% PID/DQN/SAC/MPC over a representative day with the comfort band shaded.
% Table 3: energy, comfort, cost with percentage deltas vs PID and vs MPC.
% Discuss where RL wins/loses and why (e.g., preheating using thermal mass).
\textbf{[TODO: results after experiments.]}

\subsection{Effect of the episode-reset strategy on learning}
% Fig. 3 (differentiator): learning curves under full/soft/rolling/random reset on
% the EnergyPlus FMU; compare sample efficiency and final performance. Emphasize
% this analysis is enabled by the framework's explicit reset mechanism.
\textbf{[TODO: reset ablation.]}

\subsection{Multi-rate fidelity}
% Compare a naive single-rate configuration against the native multi-rate one and
% show the difference in control quality or numerical artifacts, supporting the
% claim that native multi-rate matters.
\textbf{[TODO: multi-rate comparison.]}

\subsection{Cross-climate generalization}
% Train on one climate and test on another (e.g., Turin vs Adelaide); report
% transfer performance.
\textbf{[TODO: generalization.]}

\subsection{Discussion}
% Interpret results for practice; articulate the RL-vs-MPC trade-off (performance
% vs modeling/operational effort); argue the value of config-only realism for
% reproducible benchmarking.
\textbf{[TODO: discussion grounded in the numbers above.]}

%==============================================================================
% OPTIONAL — STAGE B2 (include to escalate to Applied Energy; otherwise remove).
% \section{Extension: prosumer co-optimization (PV and battery)}\label{sec:b2}
% Couple PV and battery to the building; the agent jointly sets heat-pump
% modulation and battery power; objective adds self-consumption and cost. Report a
% comfort/cost/self-consumption trade-off (Pareto front across reward weights).
% Stress that adding generation and storage is, again, a configuration change.
%==============================================================================

%==============================================================================
\section{Limitations and future work}\label{sec:future}
%==============================================================================
This study considers a single building controlled by a single agent. Natural
extensions are district-scale coordination with multiple buildings and shared
resources, which the underlying framework supports through multiple federations and
multi-agent learning; the use of real tariff and demand-response signals; and
sim-to-real deployment, in which the simulated building federate is replaced by a
real-time hardware federate through a configuration change while the trained policy
is reused.

%==============================================================================
\section{Conclusion}\label{sec:conclusion}
%==============================================================================
We presented an RL control study of an EnergyPlus-coupled heat-pump building in
which the realistic, multi-rate, multi-formalism environment was assembled entirely
from configuration. Because the physics was a shared, declarative artifact, we
could compare RL agents against rule-based, PID, and MPC baselines under identical
conditions and analyze how the episode-reset strategy imposed by a non-resettable
co-simulation affects learning. The results
% TODO: one-line restatement of the headline finding.
indicate that treating realistic physics as a configuration concern enables fair
and reproducible evaluation of learning-based building control, and they point
toward district-scale, multi-agent, and sim-to-real extensions.

%==============================================================================
% \section*{Acknowledgements}
% \appendix \section{Scenario configuration} full YAML + seeds.
\bibliographystyle{elsarticle-num}
% \bibliography{cosimgym}
\end{document}
